% ======================================================================
% Capítulo: Implementación
% Fases 3, 5, 6 y 7 del anteproyecto. Extensión orientativa: 8--10 páginas. Código corto solo donde aporte; el resto, a apéndices.
% ======================================================================
\chapter{Implementación}
\label{cap:implementacion}

\section{Módulo de ingesta en Go}
\label{sec:impl-go}
% TODO: Clientes de BOE y BOJA, descarga concurrente, extracción, S3, eventos y filtros por gRPC.

\section{Backend .NET: ingesta}
\label{sec:impl-ingesta}
% TODO: Consumidores, fragmentación de BOE y BOJA, municipio y fecha, embeddings, Postgres y Qdrant.

\section{Orquestación RAG}
\label{sec:impl-rag}
% TODO: Recuperación híbrida, prompts (\texttt{PromptCatalog}), citas de fuentes y streaming por SignalR.

\section{Contexto personalizado y Arena}
\label{sec:impl-arena}
% TODO: Preferencias, comparación de pipelines, votos, \emph{shadow testing}; la Arena como herramienta de validación con funcionarios.

\section{Frontend Angular}
\label{sec:impl-angular}
% TODO: Organización por features, señales, chat, visor de citas y adaptación a móvil.

\section{Orquestación y despliegue}
\label{sec:impl-despliegue}
% TODO: AppHost de Aspire, Docker Compose, Dokploy y CI/CD.

\section{Observabilidad y coste}
\label{sec:impl-observabilidad}
% TODO: OpenTelemetry, métricas GenAI (tokens, latencia, TTFT, coste en EUR), Grafana y privacidad (ADR 0002). Tabla de coste por pregunta: modelos frontera frente a locales.

\section{Auditoría y refactorización}
\label{sec:impl-auditoria}
% TODO: Auditoría de código de octubre de 2026 y sus fases de corrección.
